%============================================================================
% LICENSING NOTICE
%============================================================================
% This WIP file, upon integration into guide-v.tex, becomes part of the
% TMS/DMS/CMS Usage Guide and is released under CC BY-NC 4.0.
% For details, see LICENSE in the repository root.
%============================================================================

%============================================================================
% FALCON BMS TMS/DMS/CMS HOTAS GUIDE
% WIP FILE TEMPLATE — FINAL (Briefing)
%============================================================================

% IMPORTANTE: Este é um TEMPLATE padronizado para arquivos WIP (Work-In-Progress)
% que serão integrados ao guide.tex.

% Nomenclatura: Siga wip-naming
% Padrão: section-C{N}-S{M}[-S{K}]-{titulo}-{status}-{data}.tex
% Exemplo: section-C5-S2-cms-actuation-hotas-tables-final-2026-01-10.tex

% Status: dev | review | final | approved | deprecated
% Locação: WIP/ (ativo) | ARCHIVE/ (aprovado/descartado)

%============================================================================
% PREAMBLE COMPLETO — TMS/DMS/CMS Usage Guide for Falcon BMS 4.38.1
% Gerado: 11 February 2026
% Status: Improvements as per GitHub Issues #34 and #33
%============================================================================

\documentclass[11pt, a4paper, twoside]{report}

% --------------------------------------------------------------------------
% BASIC ENCODING AND LANGUAGE
% --------------------------------------------------------------------------

\usepackage[utf8]{inputenc}
\usepackage[T1]{fontenc}
\usepackage[english]{babel}

% --------------------------------------------------------------------------
% FONTS AND MICROTYPOGRAPHY
% --------------------------------------------------------------------------

\usepackage{lmodern}
\usepackage{microtype}

% --------------------------------------------------------------------------
% PAGE GEOMETRY AND LAYOUT
% --------------------------------------------------------------------------

\usepackage{geometry}
\geometry{a4paper, left=2.0cm, right=2.0cm, top=2.5cm, bottom=2.5cm}
\usepackage{setspace}
\onehalfspacing

% --------------------------------------------------------------------------
% COLORS AND LINKS AND PDF SETTINGS
% --------------------------------------------------------------------------

\usepackage[table]{xcolor}
\definecolor{linkblue}{HTML}{004488}
\definecolor{linkred}{HTML}{882222}
\definecolor{headerblue}{HTML}{003366}
\definecolor{rowgray}{HTML}{F5F5F5}
\definecolor{subheadgray}{HTML}{E0E0E0}
\usepackage{soul}
\usepackage[pdfencoding=auto, psdextra, colorlinks=true, linkcolor=linkblue, citecolor=linkred, urlcolor=linkblue, breaklinks=true]{hyperref}
\hypersetup{
	pdftitle={TMS, DMS and CMS Usage Guide for Falcon BMS},
	pdfauthor={Carlos "Metal" Nader},
	pdfsubject={Flight Simulation - Falcon BMS HOTAS Reference},
	pdfkeywords={Falcon BMS, F-16, HOTAS, TMS, DMS, CMS, Flight Simulation},
	pdfcreator={pdfLaTeX (MiKTeX)},
	pdfproducer={MiKTeX},
	pdflang={en-US},
}
\usepackage{bookmark}

% --------------------------------------------------------------------------
% CAPTION SETUP BY TYPE
% --------------------------------------------------------------------------

\usepackage{caption}

% GLOBAL PATTERN
\captionsetup{font=small, labelfont=bf, justification=centering, singlelinecheck=true}
% FIGURES
\captionsetup[figure]{font=footnotesize, labelfont=bf, justification=centering, singlelinecheck=true}
% % TABLE/LONGTABLE/hotastable:
\captionsetup[table]{font=small, labelfont=bf, justification=centering, singlelinecheck=true}

% --------------------------------------------------------------------------
% HEADERS AND FOOTERS (IMPROVED for report + twoside)
% --------------------------------------------------------------------------

\usepackage{fancyhdr}
\setlength{\headheight}{25pt}                    
\pagestyle{fancy}
\fancyhf{}                                        
\fancyhead[LO,RE]{\small\textit{\leftmark}}     
\fancyhead[RO,LE]{\small\thepage}               
\fancyfoot{}                                      
\renewcommand{\headrulewidth}{0.4pt}
\renewcommand{\footrulewidth}{0pt}

% --------------------------------------------------------------------------
% CHAPTER FORMATTING AND SPACING (via titlesec)
% --------------------------------------------------------------------------

\usepackage{titlesec}

% --------------------------------------------------------------------------
% CHAPTER - Nível 0 (Destaque Máximo)
% --------------------------------------------------------------------------
% Shape: display → standalone em nova linha
% Font: \Large\bfseries --- label e título ambos em bold para destaque máximo
% Numeração: SIM
% Espaçamento: 15pt antes, 28pt depois
% --------------------------------------------------------------------------

\titleformat{\chapter}[display]
{\normalfont\Large\bfseries}           % Formato: LARGE + bold (label)
{\chaptertitlename~\thechapter}        % Rótulo: "Chapter 1" (bold também)
{20pt}                                 % Espaço vertical antes do corpo
{\Large\bfseries}                      % Corpo do título: LARGE + bold

\titlespacing{\chapter}
{0pt}       % Margem esquerda (sem indent)
{15pt}      % Espaço ANTES do título do capítulo
{28pt}      % Espaço DEPOIS do título do capítulo (AUMENTADO)
[0pt]       % Margem direita

% --------------------------------------------------------------------------
% SECTION - Nível 1
% --------------------------------------------------------------------------
% Shape: hang → label à esquerda, corpo indentado (compacto)
% Font: \large\bfseries (mesmo destaque que chapter para coerência)
% Numeração: SIM
% Espaçamento: 15pt antes, 8pt depois (REDUZIDO de 15pt/10pt anteriores)
% --------------------------------------------------------------------------

\titleformat{\section}[hang]
{\normalfont\large\bfseries}           % Formato: large + bold
{\thesection}                          % Rótulo: "1", "2", etc
{1em}                                  % Espaço entre rótulo e corpo
{}                                     % Corpo do título (herda do format)

\titlespacing{\section}
{0pt}       % Margem esquerda (sem indent)
{15pt}      % Espaço ANTES do título da seção (REDUZIDO)
{8pt}       % Espaço DEPOIS do título da seção (REDUZIDO)
[0pt]       % Margem direita

% --------------------------------------------------------------------------
% SUBSECTION - Nível 2
% --------------------------------------------------------------------------
% Shape: hang → label esquerda, corpo indentado
% Font: \normalsize\bfseries (redução visual, ainda robusto)
% Numeração: SIM
% Espaçamento: 12pt antes, 6pt depois (compacto)
% --------------------------------------------------------------------------

\titleformat{\subsection}[hang]
{\normalfont\normalsize\bfseries}      % Formato: tamanho normal + bold
{\thesubsection}                       % Rótulo: "1.1", "1.2", etc
{1em}                                  % Espaço entre rótulo e corpo
{}                                     % Corpo do título

\titlespacing{\subsection}
{0pt}       % Margem esquerda (sem indent)
{12pt}      % Espaço ANTES do título da subseção
{6pt}       % Espaço DEPOIS do título da subseção
[0pt]       % Margem direita

% --------------------------------------------------------------------------
% SUBSUBSECTION - Nível 3
% --------------------------------------------------------------------------
% Shape: hang → label esquerda, máximo recuo (bem compacto)
% Font: \normalsize\bfseries (mesmo tamanho que subsection, padrão)
% Numeração: SIM (configurable com secnumdepth)
% Espaçamento: 10pt antes, 4pt depois (bem comprimido)
% --------------------------------------------------------------------------

\titleformat{\subsubsection}[hang]
{\normalfont\normalsize\bfseries}      % Formato: tamanho normal + bold
{\thesubsubsection}                    % Rótulo: "1.1.1", "1.1.2", etc
{1em}                                  % Espaço entre rótulo e corpo
{}                                     % Corpo do título

\titlespacing{\subsubsection}
{0pt}       % Margem esquerda (sem indent)
{10pt}       % Espaço ANTES do título
{4pt}       % Espaço DEPOIS do título (bem comprimido)
[0pt]       % Margem direita

% --------------------------------------------------------------------------
% PARAGRAPH - Nível 4 (CRÍTICO: Ilusão Óptica Resolvida)
% --------------------------------------------------------------------------
% Shape: runin → inline, integrado ao texto (seu requisito)
% Font: \small\bfseries (SOLUÇÃO PARA ILUSÃO ÓPTICA DO NEGRITO)%       
% Numeração: NÃO (padrão para \paragraph)
% Espaçamento: runin (sem espaço vertical antes, integrado ao texto)
% --------------------------------------------------------------------------

\titleformat{\paragraph}[runin]
{\normalfont\small\bfseries}           % Formato: SMALL + bold (compensa ilusão)
{}                                     % Rótulo: vazio (sem numeração)
{0em}                                  % Espaço entre rótulo e corpo (zero, runin)
{}                                     % Corpo do título

\titlespacing{\paragraph}
{0pt}       % Margem esquerda (sem indent)
{8pt}       % Espaço ANTES do título do parágrafo (ADICIONADO - cria separação)
{1em}       % Espaço DEPOIS do título (ADICIONADO - espaço após ":")
[0pt]       % Margem direita

% --------------------------------------------------------------------------
% SUBPARAGRAPH - Nível 5 (Preparação Futura)
% --------------------------------------------------------------------------
% Shape: runin → inline, integrado ao texto
% Font: \small (MENOS destaque que \paragraph, sem bold)
% Numeração: NÃO (padrão para \subparagraph)
% Espaçamento: runin (integrado)
% Uso: Se usado, aparecerá com destaque menor que \paragraph
% --------------------------------------------------------------------------

\titleformat{\subparagraph}[runin]
{\normalfont\small}                    % Formato: small, SEM bold (menos destaque)
{}                                     % Rótulo: vazio (sem numeração)
{0em}                                  % Espaço entre rótulo e corpo (zero, runin)
{}                                     % Corpo do título

\titlespacing{\subparagraph}
{0pt}       % Margem esquerda (sem indent)
{8pt}       % Espaço ANTES do título (ADICIONADO)
{1em}       % Espaço DEPOIS do título (ADICIONADO)
[0pt]       % Margem direita

% --------------------------------------------------------------------------
% TABLES AND MACROS
% --------------------------------------------------------------------------

\usepackage{booktabs}
\usepackage{array}
\usepackage{longtable}
\usepackage{tabularx}

% Custom Columns
\newcolumntype{L}[1]{>{\raggedright\arraybackslash}p{#1}}
\newcolumntype{C}[1]{>{\centering\arraybackslash}p{#1}}
\newcolumntype{R}[1]{>{\raggedleft\arraybackslash}p{#1}}

% Macro for Visual Reference Links
\newcommand{\imglink}[1]{\hspace{2pt}\hyperref[#1]{\scriptsize\textbf{[Fig]}}}

% --------------------------------------------------------------------------
% HOTAS LOT FILE HANDLE 
% --------------------------------------------------------------------------
\makeatletter
\newcommand{\listofhotastables}{%
	\section*{List of HOTAS Tables}%
	\addcontentsline{toc}{section}{List of HOTAS Tables}%
	\@starttoc{hotas}%
}
\makeatother

% --------------------------------------------------------------------------
% HOTAS table environment
% Version: 2.1 (2026-02-10)
% --------------------------------------------------------------------------

\newenvironment{hotastable}[1]{%
	\footnotesize
	\setlength{\tabcolsep}{3pt}
	\renewcommand{\arraystretch}{1.35}
	\begin{longtable}{L{1.00cm} L{0.90cm} L{0.90cm} L{3.30cm} L{6.40cm} L{1.40cm} L{1.60cm}}
		\caption{#1}\label{table.\thetable}\\
		\noalign{\addtocontents{hotas}{\protect\contentsline{table}{\protect\numberline{\thetable}#1}{\thepage}{table.\thetable}}}%
		\rowcolor{headerblue}
		\multicolumn{1}{>{\centering\arraybackslash}p{1.00cm}}{\textbf{\color{white}Mode}} &
		\multicolumn{1}{>{\centering\arraybackslash}p{0.90cm}}{\textbf{\color{white}Dir.}} &
		\multicolumn{1}{>{\centering\arraybackslash}p{0.90cm}}{\textbf{\color{white}Act.}} &
		\multicolumn{1}{>{\centering\arraybackslash}p{3.30cm}}{\textbf{\color{white}Function}} &
		\multicolumn{1}{>{\centering\arraybackslash}p{6.40cm}}{\textbf{\color{white}Effect / Nuance}} &
		\multicolumn{1}{>{\centering\arraybackslash}p{1.40cm}}{\textbf{\color{white}Dash34}} &
		\multicolumn{1}{>{\centering\arraybackslash}p{1.60cm}}{\textbf{\color{white}Train.}} \\
		\endfirsthead
		\rowcolor{headerblue}
		\multicolumn{1}{>{\centering\arraybackslash}p{1.00cm}}{\textbf{\color{white}Mode}} &
		\multicolumn{1}{>{\centering\arraybackslash}p{0.90cm}}{\textbf{\color{white}Dir.}} &
		\multicolumn{1}{>{\centering\arraybackslash}p{0.90cm}}{\textbf{\color{white}Act.}} &
		\multicolumn{1}{>{\centering\arraybackslash}p{3.30cm}}{\textbf{\color{white}Function}} &
		\multicolumn{1}{>{\centering\arraybackslash}p{6.40cm}}{\textbf{\color{white}Effect / Nuance}} &
		\multicolumn{1}{>{\centering\arraybackslash}p{1.40cm}}{\textbf{\color{white}Dash34}} &
		\multicolumn{1}{>{\centering\arraybackslash}p{1.60cm}}{\textbf{\color{white}Train.}} \\
		\endhead
		\multicolumn{7}{r}{\small\emph{Continued on next page}}\\
		\endfoot
		\endlastfoot
	}{%
	\end{longtable}
}

% --------------------------------------------------------------------------
% SIMPLE REFERENCE MACROS FOR BMS DOCS
% --------------------------------------------------------------------------

\newcommand{\dashref}[1]{Dash-34~\S~#1}
\newcommand{\dashone}[1]{Dash-1~\S~#1}
\newcommand{\trnref}[1]{TRN~#1}
\newcommand{\trnman}{BMS Training Manual 4.38.1}
\newcommand{\bmsver}{Falcon BMS~4.38.1}
\newcommand{\dashrefs}[1]{\textit{TO 1F-16CMAM-34-1-1}, Dash-34, sections \texttt{#1}}

% --------------------------------------------------------------------------
% CROSS-REFERENCE MACROS (INTERNAL GUIDE)
% --------------------------------------------------------------------------

\newcommand{\secref}[1]{\hyperref[#1]{Section~\ref*{#1}}}
\newcommand{\chapref}[1]{\hyperref[#1]{Chapter~\ref*{#1}}}
\newcommand{\tabref}[1]{\hyperref[#1]{Table~\ref*{#1}}}
\newcommand{\figref}[1]{\hyperref[#1]{Figure~\ref*{#1}}}

% --------------------------------------------------------------------------
% VERSION CONTROL MACROS
% --------------------------------------------------------------------------

\newcommand{\docversion}{0.4.2.0}
\newcommand{\docbuild}{19 February 2026}
\newcommand{\docstartdate}{05 January 2026}
\newcommand{\docenddate}{19 February 2026}
\newcommand{\chapterscompletedof}{0/7}
\newcommand{\tablesfilledpct}{0}
\newcommand{\fulldocversion}{\docversion+\docbuild}

% --------------------------------------------------------------------------
% GRAPHICS
% --------------------------------------------------------------------------

\usepackage{graphicx}
\graphicspath{{fig/}}
\usepackage{float}
\usepackage{enumitem}

% --------------------------------------------------------------------------
% TITLE
% --------------------------------------------------------------------------

\title{TMS, DMS and CMS Usage Guide for \bmsver}
\author{Carlos ``Metal'' Nader}
\date{Version \fulldocversion{} | Progress: Chapters \chapterscompletedof{} | Tables \tablesfilledpct{} | January 2026}

% ============================================================================
% DOCUMENT BEGIN
% ============================================================================

\begin{document}
	
	\maketitle
	
	\pagenumbering{roman}
	
	% --------------------------------------------------------------------------
	% TOC DEPTH CONFIGURATION
	% --------------------------------------------------------------------------
	\setcounter{tocdepth}{3}       % Show up to \subsubsection in TOC
	\setcounter{secnumdepth}{3}    % Number up to \subsubsection
	
	\newpage
	\tableofcontents
	\newpage
	
	% --------------------------------------------------------------------------
	% LIST OF HOTAS TABLES
	% --------------------------------------------------------------------------
	\phantomsection
	\listofhotastables
	\newpage
	
	\pagenumbering{arabic}

%============================================================================
% WIP FILE METADATA (NOT RENDERED IN PDF)
%============================================================================
% File Name: chapter-C5-style-rev-review-2026-02-24.tex
% WIP Naming Convention
% Target Chapter: C5 (CMS — Countermeasures Management Switch)
% Target Section: ALL (author review pass — full chapter)
%
% WIP Status: review
% Created: 2026-02-19
% Last Modified: 2026-02-24
%
% Narrative Completion: 100% (§5.1–§5.3 completos)
% Table Fill Status: hotastable cells revised (style pass — violations corrected)
%
% Notes:
% - Contains approved revised prose for C5 only.
% - Each passage enters this file after explicit author approval.
% - Status remains 'review' until author review pass is complete; then → final.
% - Session 2026-02-23 (author review): §5.1 approved. See issue #44 for
%   full change log.
% - Session 2026-02-24 (author review): §5.2 opening, §5.2.1 opening, §5.2.1.1
%   opening and table, §5.2.1.2 table, §5.2.1.3 table approved. Program 6
%   corrections validated across all CMDS mode tables. See issue #44 for full
%   change log. §5.2.2–§5.3 pending.
% - Session 2026-02-24 session 2 (author review): §5.2.1 opening (Program 5
%   sentence), §5.2.1.2 prose (P1, P2) and table, §5.2.1.3 prose (P1, P2)
%   and table approved. See issue #44 for full change log. §5.2.2–§5.3
%   pending.
% - PENDING: RWR introduced for the first time at §5.2.1.2. Consider conceptual
%   introduction in C2 (HOTAS Fundamentals).
% - PENDING: BYP and STBY modes not covered in §5.2.1. See sub-issue #48 (linked
%   to #44). Cascade impacts mapped: §5.2.1 opening ("three modes" → five),
%   §5.2.1.4/5 new subsubsections + hotastables, §5.2.3 hotastable Mode column,
%   §5.2.4 State Tracking, §5.1 verification.
% - PENDING: "three operational modes" in §5.2.1 opening will require update when
%   BYP/STBY are added (sub-issue #48).
% - Cross-reference: Issue #44 (v0.4.2.0 style revision / author review).

%============================================================================
% TEMPLATE STRUCTURE — Replace section/subsection headers with your content
%============================================================================

%============================================================================
% CONTENT BEGINS HERE (use \chapter{}, \section{}, \subsection{}, etc.)
%============================================================================

%============================================================================
% CHAPTER 5: CMS — Countermeasures Management Switch
%============================================================================
\chapter{CMS – Countermeasures Management Switch}
\label{chap:C5}

%============================================================================
% SECTION 5.1: Concept
%============================================================================
\section{Concept}
\label{sec:C5-S1}

The Countermeasures Management Switch (CMS) is a four-direction hat switch mounted on the side-stick controller. It controls the aircraft's defensive avionics: the ALE-47 CMDS (Countermeasures Dispenser Set — the chaff and flare dispensing system) and ECM (Electronic Countermeasures), either external pods or internal avionics, depending on block and variant.

CMS actuation procedures are detailed in \secref{sec:C5-S2}. Block and variant differences are covered in \secref{sec:C5-S3}.

For CMDS and ECM operation, see \dashrefs{2.7.1 and 2.7.2}, respectively. Diagrams of the side-stick controller and throttle are provided in \dashref{2.1.5}. \figref{fig:f16_hotas_cms_location} shows the CMS location on the F-16 side-stick controller.

\begin{figure}[H]
\centering
\includegraphics[width=0.65\textwidth]{F-16_Side_Stick_Controller-1.jpg}
\caption{F-16 Throttle and Side-Stick Controller. Image by Falconpedia (\url{falcon4.wikidot.com}), via Wikimedia Commons (\url{https://commons.wikimedia.org/wiki/File:F-16_Side_Stick_Controller.jpg}), licensed under the Creative Commons Attribution-Share Alike 3.0 Unported (CC BY-SA 3.0) license.}
\label{fig:f16_hotas_cms_location}
\end{figure}

%============================================================================
% SECTION 5.2: CMS Actuation
%============================================================================
\section{CMS Actuation}
\label{sec:C5-S2}

CMS actuation is presented here in tables organized by system: CMDS (Manual, Automatic, and Semi-Automatic modes) and ECM (external pods ALQ-131/ALQ-184 and internal IDIAS).

CMS actuation is \textit{independent} of the Master Mode currently selected.

Each table entry specifies:

\begin{itemize}
  \item Mode: The operational context (Master Mode).
  \item Direction: CMS hat direction (Up, Down, Left, Right).
  \item Action: Press type (Short, Long, Long Hold).
  \item Function: What the CMS command activates or controls.
  \item Effect / Nuance: Resulting system behavior, including tactics and constraints.
  \item Dash34: Reference section in the Dash-34 manual.
  \item Training: Recommended BMS training mission(s) (see \trnman{} for descriptions).
\end{itemize}

%============================================================================
% SECTION 5.2.1: CMDS
%============================================================================
\subsection{CMDS}
\label{sec:C5-S2-S1}

The ALE-47 CMDS operates in three modes selected via the CMDS CCU MODE knob: Manual (MAN), Automatic (AUTO), and Semi-Automatic (SEMI). The program executed in each mode --- Programs 1--4 --- is selected via the PRGM knob on the CMDS panel.

Programs 5 and 6 are independent of the PRGM knob position. Program 5 is activated via the cockpit slap switch — outside the scope of CMS — and is not covered in this guide; see \dashref{2.7.2} for details

%============================================================================
% SECTION 5.2.1.1: Manual Mode
%============================================================================
\subsubsection{Manual Mode}
\label{sec:C5-S2-S1-S1}

In Manual (MAN) mode, the pilot runs countermeasure programs via \textbf{CMS Up} and \textbf{CMS Left}.

\begin{hotastable}{CMS Actuation with CMDS Manual Mode}\label{tab:C5-S2-S1-S1-CMDS-Manual}
CMDS MAN & Up & Short & Execute Program 1--4 &
Runs the program selected via the CMDS panel PRGM knob once per press. No threat sensing; no automatic triggering. Overrides any AUTO dispensing, if running.
& 2.7.2.2 & 18, 28 \\
CMDS MAN & Left & Short & Execute Program 6 &
Program 6 can be dispensed independently of the PRGM knob position.
& 2.7.2.2 & 19 \\
\end{hotastable}

%============================================================================
% SECTION 5.2.1.2: Automatic Mode
%============================================================================
\subsubsection{Automatic Mode}
\label{sec:C5-S2-S1-S2}

In Automatic (AUTO) mode, the ALE-47 CMDS dispenses the selected program (1–4) \textit{continuously} in response to threats detected by the Radar Warning Receiver (RWR).

The pilot presses \textbf{CMS Down} \textit{once} to grant consent for AUTO dispensing; dispensing continues until the pilot presses \textbf{CMS Right} or expendables are exhausted.

\begin{hotastable}{CMS Actuation with CMDS Automatic Mode}\label{tab:C5-S2-S1-S2-CMDS-Auto}
CMDS AUTO & Up & Short & Execute Program 1--4 &
Runs the program selected via the CMDS panel PRGM knob once, interrupting any ongoing AUTO dispensing. After execution, AUTO resumes if the threat persists.
& 2.7.2.2 & 18, 28 \\
CMDS AUTO & Left & Short & Execute Program 6 &
Program 6 is dispensed independently of the PRGM knob position. Overrides AUTO; after execution, AUTO resumes.
& 2.7.2.2 & 19 \\
CMDS AUTO & Down & Short & Give Consent; Enable AUTO Dispensing &
Give Consent; Enable AUTO Dispensing &
Grants consent for AUTO dispensing. RWR-detected threats trigger automatic execution of the program selected via the CMDS panel PRGM knob. Dispensing continues until the threat clears, the pilot presses CMS Right, or expendables are exhausted. Consent state persists if the pilot switches to MAN mode; if the pilot re-engages AUTO, dispensing resumes immediately on the next detected threat.
& 2.7.2.1 & 18, 28 \\
CMDS AUTO & Right & Short & Cancel Consent; Disable AUTO &
Removes CMDS consent and places the ALE-47 in Standby. Automatic dispensing halts immediately. The pilot must re-issue CMS Down to resume AUTO operation.
& 2.7.2.1 & 18, 28 \\
\end{hotastable}

%============================================================================
% SECTION 5.2.1.3: Semi-Automatic Mode
%============================================================================
\subsubsection{Semi-Automatic Mode}
\label{sec:C5-S2-S1-S3}

In Semi-Automatic (SEMI) mode, the ALE-47 CMDS dispenses \textit{one program} per detected threat; the pilot must grant consent before dispensing occurs.

When a threat is detected, the CMDS control unit displays \texttt{DISPENSE RDY} and sounds a \texttt{COUNTER} voice message. If the threat persists or a new threat appears, the CMDS issues \texttt{COUNTER} again.

\begin{hotastable}{CMS Actuation with CMDS Semi-Automatic Mode}\label{tab:C5-S2-S1-S3-CMDS-Semi}
CMDS SEMI & Up & Short & Execute Program 1--4 &
Runs the program selected via the CMDS panel PRGM knob once, independent of RWR threat state. After execution, the CMDS returns to standby and resumes threat detection.
& 2.7.2.2 & 18, 28 \\
CMDS SEMI & Left & Short & Execute Program 6 &
Program 6 is dispensed independently of the PRGM knob position. After execution, the CMDS returns to standby and resumes threat detection.
& 2.7.2.2 & 19 \\
CMDS SEMI & Down & Short & Give Consent; Dispense One Program &
Dispenses one program in response to a \texttt{COUNTER} prompt. If the threat persists or a new threat appears, the CMDS issues \texttt{COUNTER} again. Consent state is retained; if the pilot switches to AUTO, the CMDS begins dispensing immediately on the next detected threat.
& 2.7.2.1 & 18, 28 \\
CMDS SEMI & Right & Short & Cancel Consent; Return to Standby &
Removes SEMI consent and places the CMDS in Standby. \texttt{COUNTER} messages cease.
& 2.7.2.1 & 18, 28 \\
\end{hotastable}

%============================================================================
% SECTION 5.2.2: ECM
%============================================================================
\subsection{ECM}
\label{sec:C5-S2-S2}

External ECM pods (ALQ-131, ALQ-184) and internal IDIAS (Improved Defensive Internal Avionic System) provide active jamming across radar frequency bands. \textbf{CMS Down} grants consent to external pods. \textbf{CMS Left} cycles operational modes for internal IDIAS. 

%============================================================================
% SECTION 5.2.2.1: ECM Pod
%============================================================================
\subsubsection{External ECM Pod}
\label{sec:C5-S2-S2-S1}

\textbf{CMS Down} grants consent for ECM pod transmission. The XMIT switch on the ECM control panel sets the transmission mode: XMIT 1 (Avionics Priority, AFT antenna only), XMIT 2 (ECM Priority, both FWD and AFT antennas), or XMIT 3 (Active Jam, continuous transmission). The ECM Enable light on the miscellaneous panel reflects consent state. The ECM system is \textit{inhibited} when the pilot moves the RF switch from NORM, or when the landing gear handle is down.

\begin{hotastable}{CMS Actuation with External ECM Pod}\label{tab:C5-S2-S2-S1-ECM-Pod}
ECM Pod & Down & Short & Grant ECM Consent &
CMS Down illuminates the ECM Enable light. The pod transmits in the mode set by the XMIT switch. Transmission continues until CMS Right is pressed or the RF switch is moved from \texttt{NORM}.
& 2.7.4.1.1, 2.7.4.2.5 & 28 \\
ECM Pod & Right & Short & Remove ECM Consent &
CMS Right removes consent and places the pod in Standby. The ECM Enable light extinguishes. Transmission does not resume until the pilot re-issues CMS Down.
& 2.7.4.1.1, 2.7.4.2.5 & 28 \\
\end{hotastable}

%============================================================================
% SECTION 5.2.2.2: Internal ECM (IDIAS)
%============================================================================
\subsubsection{Internal ECM (IDIAS)}
\label{sec:C5-S2-S2-S2}

IDIAS provides three operational modes controlled via \textbf{CMS Left}: Standby (\texttt{STBY}), Avionics Priority (\texttt{AVNC}), and ECM Priority (\texttt{ECM}). Operational modes are available only when the XMTR button is set to OPER. The active mode is reflected on the HUD, FCR, RWR, and TFR pages.

\begin{hotastable}{CMS Actuation with Internal ECM (IDIAS)}\label{tab:C5-S2-S2-S2-IDIAS}
IDIAS ECM & Left & Short & Cycle ECM Operational Mode &
Each short press of CMS Left advances the IDIAS mode: \texttt{STBY} $\rightarrow$ \texttt{AVNC} (Avionics Priority) $\rightarrow$ \texttt{ECM} (ECM Priority) $\rightarrow$ \texttt{AVNC} $\rightarrow$ \texttt{ECM}.
& 2.7.4.1.2 & 28 \\
IDIAS ECM & Right & Short & Remove ECM Consent &
CMS Right sets IDIAS to \texttt{STBY}, removing consent. \texttt{AVNC} or \texttt{ECM} modes are not restored until CMS Left is re-applied.
& 2.7.4.1.2 & 28 \\
\end{hotastable}

%============================================================================
% SECTION 5.2.3: Consent and Constraints
%============================================================================
\subsection{Consent and Constraints}
\label{sec:C5-S2-S3}

When equipped with an external ECM pod, \textbf{CMS Down} grants consent to both CMDS (in AUTO or SEMI mode) and the external ECM pod \textit{simultaneously}.

\begin{hotastable}{CMS Interaction with CMDS and ECM (Consent and Constraints)}\label{tab:C5-S2-S3-CMDS-ECM}
CMDS AUTO SEMI + ECM Pod & Down & Short & Joint Consent (CMDS + ECM) &
A \textit{single} CMS Down press grants consent to \textit{both} the ALE-47 CMDS (AUTO or SEMI mode) and the external ECM pod.
& 2.7.2.1, 2.7.4.2.5 & 18, 28 \\
\end{hotastable}

%============================================================================
% SECTION 5.2.4: Operational Notes
%============================================================================
\subsection{Operational Notes}
\label{sec:C5-S2-S4}

\paragraph{State Tracking:}

In AUTO and SEMI modes, the CMDS retains consent state when the pilot switches to MAN mode. If the pilot re-engages AUTO without pressing \textbf{CMS Down}, the CMDS begins dispensing immediately on the next detected threat.

\paragraph{Bingo Quantity Behavior:}

If expendables fall to or below the bingo quantity, the CMDS requests consent and continues dispensing. The VMU issues a \texttt{LOW} message when an expendable reaches the bingo quantity, and an \texttt{OUT} message when an expendable is depleted. Dispensing does not stop automatically. The pilot can press \textbf{CMS Right} to inhibit AUTO dispensing, or switch to MAN mode to control expenditure directly.

\paragraph{RF Switch Override:}

When ECM consent is active, moving the RF switch from NORM to QUIET or SILENT places the ECM system in Standby. Returning the RF switch to NORM does \textit{not} restore consent. The pilot must re-issue \textbf{CMS Down} to resume ECM transmission.

\paragraph{IDIAS vs. External ECM Pod:}

The joint consent described in \tabref{tab:C5-S2-S3-CMDS-ECM} applies to external ECM pod configurations \textit{only}. On IDIAS-equipped aircraft, \textbf{CMS Left} is used to cycle ECM modes; \textbf{CMS Down} has no ECM function on these variants.

\paragraph{Ground Safety:}

On the ground, ECM consent state is Standby. If the pilot holds \textbf{CMS Down} on the ground, the ECM pod radiates until the pilot releases \textbf{CMS Down}. Ground personnel must be clear of the radiation area if, for any reason, ECM must be enabled on the ground.

%============================================================================
% SECTION 5.3: ECM Configuration by Block and Variant
%============================================================================
\section{ECM Configuration by Block and Variant}
\label{sec:C5-S3}

CMS interaction with CMDS is \textit{uniform} across all F-16 blocks and variants (see \secref{sec:C5-S2-S1}). External Pod and IDIAS F-16 variants, however, use \textit{different} CMS procedures and ECM control panels (see \secref{sec:C5-S2-S2}):

\begin{table}[H]
\centering
\caption{ECM Configuration by Variant --- CMS Interface Summary}
\label{tab:C5-S3-ecm-config-summary}
\small
\renewcommand{\arraystretch}{1.25}
\begin{tabular}{|>{\raggedright\arraybackslash}p{3.6cm}|>{\raggedright\arraybackslash}p{3.6cm}|>{\raggedright\arraybackslash}p{4.6cm}|}
\hline
\rowcolor{headerblue}
\textcolor{white}{\textbf{ECM Configuration}} & \textcolor{white}{\textbf{CMS Direction}} & \textcolor{white}{\textbf{Panel Control}} \\
\hline
External Pod (ALQ-131/ALQ-184) & CMS Down (consent) & XMIT knob (positions: 1, 2, 3) \\
\hline
Internal IDIAS & CMS Left (mode cycling) & XMTR switch (positions: STBY, OPER) \\
\hline
\end{tabular}
\end{table}

%============================================================================
% SECTION 5.3.1: ECM Configurations in BMS
%============================================================================
\subsection{ECM Configurations in BMS}
\label{sec:C5-S3-S1}

Falcon BMS 4.38.1 implements two distinct ECM system architectures, each with different CMS actuation methods. Detailed actuation procedures are provided in \secref{sec:C5-S2}.

The variants listed in the following tables reflect Dash-34 documentation; not all may be available as selectable aircraft in Falcon BMS 4.38.1.

%============================================================================
% SECTION 5.3.1.1: ECM Pods (ALQ-131 / ALQ-184)
%============================================================================
\subsubsection{ECM Pods (ALQ-131 / ALQ-184)}
\label{sec:C5-S3-S1-S1}

These variants use the ECM Pod Control Panel with manual jamming band selection. \textbf{CMS Down} grants consent for ECM transmission. For detailed procedures, see \secref{sec:C5-S2-S2-S1}.

\begin{table}[H]
\centering
\caption{External ECM Pod Variants}
\label{tab:C5-S3-external-ecm-pods}
\small
\renewcommand{\arraystretch}{1.25}
\begin{tabular}{|>{\raggedright\arraybackslash}p{3.0cm}|>{\raggedright\arraybackslash}p{5.8cm}|}
\hline
\rowcolor{headerblue}
\textcolor{white}{\textbf{Operator}} & \textcolor{white}{\textbf{Block/Variant}} \\
\hline
USAF & Blocks 40/42/50/52 (CM designation) \\
\hline
NATO & Block 15 operators (Belgium, Denmark, Netherlands, Norway) \\
\hline
International & Egypt, Korea KF-16C Block 32 \\
\hline
\end{tabular}
\end{table}

%============================================================================
% SECTION 5.3.1.2: Internal ECM (IDIAS)
%============================================================================
\subsubsection{Internal ECM (IDIAS)}
\label{sec:C5-S3-S1-S2}

These variants use the IDIAS ECM control panel. \textbf{CMS Left} cycles the operational mode. For detailed procedures, see \secref{sec:C5-S2-S2-S2}.

\begin{table}[H]
\centering
\caption{Internal ECM (IDIAS) Variants}
\label{tab:C5-S3-internal-ecm}
\small
\renewcommand{\arraystretch}{1.25}
\begin{tabular}{|>{\raggedright\arraybackslash}p{3.0cm}|>{\raggedright\arraybackslash}p{5.8cm}|}
\hline
\rowcolor{headerblue}
\textcolor{white}{\textbf{Operator}} & \textcolor{white}{\textbf{Block/Variant}} \\
\hline
Israel & F-16I Sufa Block 52, Barak I Block 30, Barak II Block 40 \\
\hline
International & Greek HAF (Blocks 50 PXII, 52 PXIII, 52+ PXIV Advanced),
Korea KF-16C Block 52, Singapore F-16D Block 52 \\
\hline
\end{tabular}
\end{table}

%============================================================================
% END OF SECTION
%============================================================================

\end{document}